\chapter{系统实现与关键模块}

\section{关键实现文件与模块关系}
本系统采用“统一平台入口调度、算法模块独立实现”的代码组织方式。平台入口文件为 \texttt{src/app.py}。该文件负责 Gradio 界面构建、用户参数接收、默认路径识别、模型权重加载和外部脚本调度。MRI 数据预处理由 \texttt{src/preprocess\_mri.py} 实现。二维生成训练与推理复用 StyleGAN2-ADA 工程。点云初始化由独立脚本完成。3DGS 相关代码和训练产物分别保存在三维框架目录和结果目录中。这样的划分使界面层、二维生成模块和三维重建模块保持相对独立。

\begin{table}[H]
\centering
\caption{系统关键实现文件及功能说明}
\begin{tabular}{>{\raggedright\arraybackslash}p{4cm}p{8.5cm}}
\toprule
文件或目录 & 主要作用 \\
\midrule
\thesispath{src/app.py} & 构建 Gradio 平台界面，调度预处理、训练、生成、序列构建、点云初始化和 AI 助手功能。 \\
\thesispath{src/preprocess_mri.py} & 读取 NIfTI 格式 MRI 体数据，提取轴向切片，完成归一化、缩放、低质量切片过滤和 PNG 保存。 \\
\thesispath{src/stylegan2} & 提供 StyleGAN2-ADA 的数据打包、训练、生成器加载和图像推理能力。 \\
\thesispath{src/init_pcd.py} & 将连续二维切片序列转换为 ASCII PLY 点云，作为 3DGS 三维表达的输入准备。 \\
\thesispath{src/gaussian_splatting} & 保存 3D Gaussian Splatting 框架代码，支撑 MRI 切片序列的三维高斯训练、渲染和点云导出。 \\
\thesispath{results/training_runs} & 保存训练配置、训练日志、真实样本网格和模型权重快照。 \\
\thesispath{results/generated_slices} & 保存平台根据训练权重生成的 MRI 切片结果。 \\
\thesispath{results/3dgs_runs} & 保存 3DGS 训练配置、相机参数、checkpoint、渲染图像和训练后点云结果。 \\
\bottomrule
\end{tabular}
\end{table}

\section{MRI 数据预处理模块实现}
MRI 数据预处理模块由 \texttt{preprocess\_mri.py} 实现。脚本接收输入目录、输出目录和目标尺寸等参数。脚本使用 \texttt{nibabel} 读取 \texttt{.nii} 或 \texttt{.nii.gz} 格式的 NIfTI 文件。脑部 MRI 体数据两端常包含较大背景或低信息区域。因此，系统默认沿轴向提取中间约 15\% 到 85\% 的切片。这样可以减少头顶部、颈部和空白切片对训练集的影响。

灰度处理阶段，系统采用 2\% 到 98\% 百分位截断。系统不直接使用全局最小值和最大值进行拉伸。这样可以降低异常亮点、噪声和局部极值的影响。归一化后的切片被转换为 \texttt{uint8} 类型。系统再通过双三次插值把切片缩放到目标分辨率。为了控制训练集质量，脚本统计切片平均灰度、标准差和非零像素比例。只有满足信息量要求的切片才会保存。

预处理完成后，平台可以调用 StyleGAN2 的 \texttt{dataset\_tool.py}。该工具把 PNG 切片目录打包为训练数据。本次训练使用的数据集文件为 \texttt{mri128.zip}。训练集最大规模为 60815 张。图像分辨率为 $128\times128$。该流程实现了从医学体数据到二维生成模型训练样本的转换。

\section{StyleGAN2 训练与切片生成实现}
StyleGAN2 训练由平台入口 \texttt{app.py} 封装。平台会在训练数据目录中查找 zip 数据包。平台优先使用本项目配置的 \texttt{mri128.zip}。训练任务启动后，系统通过 \texttt{subprocess.Popen} 调用 StyleGAN2 训练脚本。系统向子进程传入输出目录、数据路径、GPU 数量、训练 kimg、batch size、学习率和快照保存频率等参数。训练耗时较长，所以任务在后台运行。这样可以避免 Gradio 主界面被阻塞。

切片生成模块使用训练权重中的 \texttt{G\_ema} 生成器进行推理。平台可以自动查找最新网络快照。用户也可以手动指定权重路径。推理阶段，系统根据随机种子生成潜变量。系统再结合截断参数 \texttt{truncation\_psi} 调节输出的稳定性和多样性。生成器输出张量从 $[-1,1]$ 转换到 $[0,255]$ 后保存为 PNG 图像。本项目已经形成编号为 010000 的模型快照，并生成了示例切片。

该模块的重点是把训练产物接入平台。用户可以通过界面设置生成数量、随机种子和截断参数。系统负责保存图像、返回路径并展示结果。该功能方便论文插图整理，也便于答辩时说明模型输入、参数和输出之间的关系。

\section{连续切片序列与点云初始化实现}
连续切片序列由 \texttt{app.py} 中的 \texttt{generate\_sequence} 函数生成。系统分别根据两个随机种子构造潜变量 $z_a$ 和 $z_b$。系统再对二者进行线性插值，得到一组连续变化的潜变量。每个插值潜变量输入 StyleGAN2 生成器后输出一张 MRI 切片。系统按 \texttt{slice\_0000.png}、\texttt{slice\_0001.png} 等规则保存图像。该功能为 2D/3D 对比展示提供连续输入。

点云初始化由 \texttt{init\_pcd.py} 完成。脚本读取连续切片目录中的 PNG 图像。脚本把灰度值大于阈值 10 的像素转换为空间点。对于第 $z$ 张切片中的像素坐标 $(x,y)$，系统将其映射为三维点 $(x,y,z\cdot spacing)$。其中，\texttt{spacing} 表示切片间距。像素灰度值归一化后被写入 RGB 颜色字段。所有点最终以 ASCII PLY 格式保存。本项目已生成 \texttt{results/3d\_model\_test/init\_point\_cloud.ply}。该文件包含 477426 个顶点，大小约 13 MB。

通过该模块，二维切片被组织为三维数据。每个点都有空间坐标和灰度颜色。该点云可以作为 3DGS 训练输入。该点云也可以展示二维生成结果向三维表达过渡的过程。

\section{3DGS 三维重建模块实现}
项目保留了 3DGS 框架目录。系统围绕 MRI 切片序列、初始点云、相机参数、训练日志和输出点云设计统一接口。该实现路线对应本文的工程目标。系统先把二维生成结果组织为三维输入，再通过 3DGS 优化形成高斯表达。

3DGS 模块输入包括三类数据。第一类是连续 MRI 切片序列。它们用于提供渲染监督。第二类是根据切片灰度和空间位置初始化得到的 PLY 点云。第三类是描述切片投影关系的相机参数文件。普通 3DGS 场景通常借助 COLMAP 估计相机位姿。MRI 切片序列具有规则层间关系。因此，系统将切片索引、像素间距和层间距转换为规则相机参数。当前训练目录中的 \texttt{cameras.json} 包含 180 个相机条目。图像分辨率为 $256\times256$。焦距参数 $fx$ 与 $fy$ 均为 426.67。

训练实现上，系统使用 180 张连续 MRI 切片和 300000 个初始化点构建 3DGS 输入数据。系统把结果保存到 3DGS 主训练目录。训练过程记录相机参数、配置参数、训练日志、checkpoint、渲染图像、参考切片和迭代点云。第 30000 轮输出的 PLY 文件头部显示包含 572073 个高斯点。系统同时保存 \texttt{chkpnt30000.pth}，用于记录训练后的模型状态。

从平台角度看，三维重建模块接收切片目录和输出目录等核心参数。模块返回训练日志、点云路径和渲染结果路径。系统先完成点云初始化，再组织训练迭代次数、输出目录、相机参数和渲染结果目录等信息。平台沿用日志框和三维展示区域。用户可以在同一界面中查看、下载和归档三维结果。

本文将该部分定义为 3DGS 三维重建模块。原因是项目已经形成真实切片序列、PLY 点云文件、训练日志、checkpoint 和渲染截图。该模块完成了从二维切片组织到三维高斯表达结果的工程闭环。第 6 章将进一步分析相关训练指标和渲染效果。

\section{平台前端与交互功能实现}
平台前端基于 Gradio Blocks 构建。平台通过多标签页组织核心功能。数据预处理页用于设置原始数据目录、输出目录和目标尺寸。模型训练页用于设置学习率、批大小和训练 kimg。医学切片生成页用于加载权重并批量生成 MRI 切片。2D/3D 对比页用于生成连续切片序列、启动点云初始化并展示三维输出。

工程实现中，平台使用当前运行环境中的 Python 解释器启动外部脚本。这样可以减少界面进程和子进程依赖不一致的问题。训练任务使用后台进程。预处理、推理生成和点云初始化等较短任务直接返回日志、图像或输出路径。平台还封装了 \texttt{\_find\_latest\_pkl} 和 \texttt{\_find\_training\_data\_zip} 等辅助函数。这些函数用于自动定位最新权重和训练数据。

\section{AI 助手功能实现}
AI 助手以悬浮按钮和聊天面板形式集成在平台右下角。系统设置了面向 MRI 医学图像生成平台的提示词。助手回答时主要围绕 StyleGAN2、3DGS、MRI 图像处理和平台使用流程展开。平台通过兼容 OpenAI SDK 的调用方式接入外部大模型服务。平台也支持在不同模型提供商之间切换。

AI 助手的定位是教学辅助和参数解释。它不是核心算法依赖。用户可以询问截断参数、随机种子、点云初始化、3DGS 三维表达和 MRI 切片生成等问题。AI 助手依赖外部接口。它的可用性会受到网络和 API 服务影响。因此，系统的数据处理、模型训练、图像生成和点云初始化功能均保持独立运行。即使 AI 助手不可用，平台的主要毕业设计功能仍能正常完成。
